% ======================================================================
%  FIX 1.3 — Corrected coarea/substitution section for lecA1.tex
%
%  This file contains two replacement blocks:
%    (A) Corrected introductory paragraphs  (replaces lines 7--33)
%    (B) Corrected coarea section           (replaces lines 137--191)
%
%  Changes address challenges:
%    ch-f87192e5  Dimension condition  m <= n  -->  m < n  (strict)
%    ch-74d9cf11  1D corollary re-attributed to area formula / substitution
%    ch-fc52e254  Honest framing: proof uses substitution, not coarea
%    ch-5ae5f061  Hausdorff = Lebesgue note added
%    ch-8f1f7e36  Delta composition gap filled
%    ch-abbda19d  (node 1.1) dimension condition now correct
%    ch-06bc298b  (node 1.1) 1D specialisation now correctly sourced
% ======================================================================

% ~~~~~~~~~~~~~~~~~~~~~~~~~~~~~~~~~~~~~~~~~~~~~~~~~~~~~~~~~~~~~~~~~~~~~~
%  (A)  INTRODUCTORY PARAGRAPHS  —  replaces lines 7--33
% ~~~~~~~~~~~~~~~~~~~~~~~~~~~~~~~~~~~~~~~~~~~~~~~~~~~~~~~~~~~~~~~~~~~~~~

\section{Addendum: the energy-momentum tensor from microscopics}%
\label{sec:emt-rigorous}

In Lecture~6 (\S\ref{sec:energy-momentum}) we defined the
energy-momentum tensor of a system of point particles via the
microscopic expressions~\eqref{eq:Tmunu-dt}
and~\eqref{eq:Tmunu-covariant}, and left the proof that these two
forms are equivalent---and that the result is a Lorentz tensor---as
an exercise (Exercise~\ref{ex:Tmunu-tensor}).  The standard argument,
due to Weinberg~\cite{Weinberg72}, consists of inserting a
coordinate-time integral, observing that ``the differentials~$dt'$
cancel,'' and replacing~$dt'$ with the invariant proper-time
element~$d\tau$.  While physically convincing, this argument glosses
over the distributional nature of the objects involved and leaves the
equivalence of the two forms essentially unproved.

In this addendum we supply the missing rigour.  For the equivalence
theorem itself (\S\ref{sec:equivalence}), the key step is the
classical \textbf{change-of-variables (substitution) formula} for a
$C^1$ diffeomorphism---an elementary result that needs no heavy
machinery.  However, the natural \emph{generalisation} to slicing by
arbitrary Cauchy surfaces (\S\ref{sec:general-slicing}) requires the
\textbf{coarea formula} of geometric measure
theory~\cite{Federer69}, which provides the framework for
decomposing integrals over level sets of Lipschitz maps.
We therefore state the coarea formula in its general form and explain
how the one-dimensional case used in the equivalence proof reduces to
the classical substitution rule, while the full coarea/slicing
machinery becomes essential for the general slicing theorem.

% ~~~~~~~~~~~~~~~~~~~~~~~~~~~~~~~~~~~~~~~~~~~~~~~~~~~~~~~~~~~~~~~~~~~~~~
%  (B)  COAREA SECTION  —  replaces lines 137--191
% ~~~~~~~~~~~~~~~~~~~~~~~~~~~~~~~~~~~~~~~~~~~~~~~~~~~~~~~~~~~~~~~~~~~~~~

\subsection{The coarea formula and the substitution rule}%
\label{sec:coarea}

The bridge between the covariant
form~\eqref{eq:Tmunu-delta4-add} and the $3{+}1$
splitting~\eqref{eq:Tmunu-dt} is ultimately the classical
\emph{change-of-variables formula} for a diffeomorphism.  To place
this in its proper context---and to prepare for the general slicing
theorem of \S\ref{sec:general-slicing}---we state the coarea formula
of Federer and then discuss the one-dimensional case separately.

\begin{theorem}[Coarea formula {\cite[Theorem~3.2.12]{Federer69}}]%
\label{thm:coarea}
  Let $f\colon \Rn{n} \to \Rn{m}$ be a Lipschitz map with
  $m < n$ (strict inequality).  Then for any nonnegative Borel function
  $g\colon \Rn{n} \to [0,\infty]$,
  \begin{equation}\label{eq:coarea-general}
    \int_{\Rn{n}} g(x)\,J_m f(x)\;\dd\Leb^n(x)
      = \int_{\Rn{m}}
        \biggl(\int_{f^{-1}(y)}
          g\;\dd\mathcal{H}^{n-m}\biggr)
        \dd\Leb^m(y)\,,
  \end{equation}
  where $J_m f$ denotes the $m$-dimensional Jacobian of~$f$,
  $\Leb^k$ is the $k$-dimensional Lebesgue measure, and
  $\mathcal{H}^{n-m}$ is the $(n{-}m)$-dimensional Hausdorff
  measure on the level set~$f^{-1}(y)$.
\end{theorem}

\begin{remark}\label{rem:Hausdorff-Lebesgue}
  On $\Rn{k}$, the $k$-dimensional Hausdorff measure~$\mathcal{H}^k$
  coincides with the Lebesgue measure~$\Leb^k$
  \cite[2.10.35]{Federer69}.  We use~$\Leb^k$ for the ambient
  measures on domain and codomain and reserve~$\mathcal{H}^{n-m}$ for
  the measure on the (typically lower-dimensional) level sets.
\end{remark}

\noindent
\textbf{The one-dimensional case ($n = m = 1$).}  When domain and
codomain have the same dimension, the strict inequality $m < n$ is
\emph{not} satisfied, so Theorem~\ref{thm:coarea} does not apply.
The relevant result is instead the \textbf{area formula}
\cite[Theorem~3.2.3]{Federer69}: for $f\colon \R \to \R$ Lipschitz
and $g\colon \R \to [0,\infty]$ Borel,
\begin{equation}\label{eq:area-1d}
  \eqbox{
  \int_{\R} g(\tau)\,|f'(\tau)|\;\dd\tau
    = \int_{\R}
      \sum_{\tau\in f^{-1}(t)} g(\tau)\;\dd t}\,.
\end{equation}
When $f$ is injective (in particular, a diffeomorphism), each level
set $f^{-1}(t)$ is at most a singleton and the sum is trivial: one
recovers the classical \textbf{substitution (change-of-variables)
formula}
\begin{equation}\label{eq:substitution}
  \int_{\R} g(\tau)\,|f'(\tau)|\;\dd\tau
    = \int_{\R} g\bigl(\tau_\ast(t)\bigr)\;\dd t\,,
\end{equation}
where $\tau_\ast(t) = f^{-1}(t)$ is the unique preimage.  This is
the identity that drives the proof of
Theorem~\ref{thm:equivalence}; the full coarea machinery is not
needed for that argument but becomes essential in
\S\ref{sec:general-slicing}.

\begin{remark}[Delta-function composition rule]%
\label{rem:delta-composition}
  The substitution formula~\eqref{eq:substitution} is equivalent to
  the distributional identity
  \begin{equation}\label{eq:delta-composition}
    \delta\bigl(f(\tau) - t\bigr)
      = \frac{\delta(\tau - \tau_\ast)}{|f'(\tau_\ast)|}\,.
  \end{equation}
  To see this, note that \eqref{eq:substitution} states that
  the pushforward of the measure $|f'(\tau)|\,\dd\tau$ under~$f$
  equals Lebesgue measure~$\dd t$ (by the Riesz representation
  theorem, since the identity holds for all nonneg.\ Borel~$g$).
  Testing against a Dirac mass at~$t$ then
  yields~\eqref{eq:delta-composition}.  More generally, when $f$ has
  finitely many roots $\{\tau_k\}$ of $f(\tau) = t$, the area
  formula~\eqref{eq:area-1d} gives the standard multi-root identity
  $\delta(f(\tau) - t) = \sum_k
  \delta(\tau - \tau_k)/|f'(\tau_k)|$.  The substitution formula thus
  provides the rigorous foundation for the ``delta-function
  composition rule'' that is often invoked without proof in physics
  texts.
\end{remark}
